% TOP_PROBLEM: {"id":30007733,"problem_number":"LOCAL-30007733","title":"Home-school investment complementarities","category_id":21,"category":{"id":21,"name":"economics","display_name":"Economics","description":"Open economic theory statements, published research agendas, and proposed scoped specifications; question provenance and historical priority are qualified per record.","slug":"economics","order_index":21,"created_at":"2026-10-08T00:00:00Z"},"status":"research_question","external_url":null,"legacy_ids":[],"published":false,"statement":"Use factorial offers to test whether one preschool curriculum and one parenting program complement each other in school readiness and five-year family welfare.","economics":{"original_id":222,"field":"Education, families and demography","kind":"empirical","formalization_basis":"adapted_research_question","source_status":"research_agenda","priority_status":"earliest_found","priority_note":"Earliest explicit statement located in this audit; no exhaustive priority search or first-ever claim. The mathematical specification below is adapted, not quoted from the source.","earliest_verified_reference":{"authors":["Greg J. Duncan","Ariel Kalil","Magne Mogstad","Mari Rege"],"title":"Investing in Early Childhood Development in Preschool and at Home","year":2022,"url":"https://bfi.uchicago.edu/wp-content/uploads/2022/05/BFI_WP_2022-58.pdf","doi":"10.3386/w29985","locator":"§7 Knowledge needs, printed pp. 91–92","evidence_summary":"The authors ask for long-term curriculum evidence and manipulation of follow-on school investments and combined preschool and parenting interventions."},"current_reference":{"authors":["Greg J. Duncan","Ariel Kalil","Magne Mogstad","Mari Rege"],"title":"Investing in Early Childhood Development in Preschool and at Home","year":2022,"url":"https://bfi.uchicago.edu/wp-content/uploads/2022/05/BFI_WP_2022-58.pdf","doi":"10.3386/w29985","locator":"§7 Knowledge needs, printed pp. 91–92","evidence_summary":"The authors ask for long-term curriculum evidence and manipulation of follow-on school investments and combined preschool and parenting interventions."},"formalization_note":"The source explicitly identifies the underlying gap or agenda. Population, horizon, interventions, estimands, and auxiliary assumptions are proposed operational choices, not a canonical source theorem. Partial later answers are compatible with this scoped question; the citation alone does not prove absence of every subsequent solution."},"statusQualification":"Published question or research agenda verified at the cited locator. The mathematical specification is adapted for this catalogue and includes additional assumptions; source publication does not establish first-ever priority or certify this exact model remains unsolved. The source explicitly identifies the underlying gap or agenda. Population, horizon, interventions, estimands, and auxiliary assumptions are proposed operational choices, not a canonical source theorem. Partial later answers are compatible with this scoped question; the citation alone does not prove absence of every subsequent solution.","statusReviewedAt":"2026-10-08"}
% FIRST_POSTED: 2026-10-08
% ENTRY_KIND: statement_only
% !TeX program = lualatex
\documentclass[11pt]{article}
\usepackage{fontspec}
\usepackage[a4paper,margin=25mm]{geometry}
\usepackage{amsmath,amssymb,mathtools}
\usepackage{xurl}
\usepackage[hidelinks]{hyperref}
\newcommand{\sref}[2]{\hyperlink{source:#1:#2}{[S#2]}}
\setlength{\emergencystretch}{3em}
\begin{document}
% problemId:30007733
\section*{Home-school investment complementarities}\label{30007733}
\subsection{Definitions and mathematical statement}
Consider four-year-old children in disadvantaged households in one urban school district. Let \(p\in\{0,1\}\) be a randomized offer of a specified preschool curriculum and \(h\in\{0,1\}\) a randomized offer of a home-visiting parenting program. Define \(S_i(p,h)\) as independently assessed school readiness one year later, and \(V_i(p,h)\) as discounted family welfare over five years net of participation time and program resource costs. The interaction targets are
\[\gamma_S=E[S_i(1,1)-S_i(1,0)-S_i(0,1)+S_i(0,0)],\]
and \(\gamma_V\), defined by replacing readiness with welfare. Identify both quantities using factorial offers, with consistency and either negligible household spillovers or a stated neighborhood exposure design. Positive \(\gamma_S\) establishes complementarity of these offers on this outcome scale, not necessarily technological complementarity of realized educational inputs. Collect attendance, parental time, private expenditure, and curriculum implementation to evaluate that distinction. Report interaction uncertainty and baseline heterogeneity instead of inferring complementarity from separate positive average effects. The review explicitly calls for combined preschool and parenting evidence and studying parental choices. Its general agenda is adapted here to defined programs and welfare accounting; no sign restriction on either interaction is assumed.

\par\textbf{Formulation and scope.} The source explicitly identifies the underlying gap or agenda. Population, horizon, interventions, estimands, and auxiliary assumptions are proposed operational choices, not a canonical source theorem. Partial later answers are compatible with this scoped question; the citation alone does not prove absence of every subsequent solution.

\par\textbf{Open-question provenance.} An explicit question or research agenda was verified at the source locator. This does not by itself certify that every later version remains unresolved \sref{30007733}{1}.

\par\textbf{Historical priority.} Earliest explicit statement located in this audit; no exhaustive priority search or first-ever claim. The mathematical specification below is adapted, not quoted from the source.

\subsection{Short English statement}
Use factorial offers to test whether one preschool curriculum and one parenting program complement each other in school readiness and five-year family welfare.

\subsection{Sources}
\begin{itemize}\raggedright
\item[\textnormal{[S1]}]\hypertarget{source:30007733:1}{} Greg J. Duncan, Ariel Kalil, Magne Mogstad, Mari Rege. 2022. Investing in Early Childhood Development in Preschool and at Home. §7 Knowledge needs, printed pp. 91–92.
\url{https://bfi.uchicago.edu/wp-content/uploads/2022/05/BFI_WP_2022-58.pdf}.
DOI: 10.3386/w29985.
{\small\textit{Earliest verified question/agenda reference; historical priority qualified below. The authors ask for long-term curriculum evidence and manipulation of follow-on school investments and combined preschool and parenting interventions.}}
\end{itemize}
\end{document}
